%============================================================
% constraints.tex — from the constraints to the algorithm (one-page guide)
% Grow the second table as chapters are written.
%============================================================
\chapter{From Constraints to Algorithm}
\label{ch:constraints}

Bit problems announce themselves in their limits. Read $n$ (how many numbers), $V$ (how large they are, often $2^{20}$, $2^{30}$ or $2^{60}$) and $q$ (queries), and look the budget up here. A judge runs roughly $10^8$ simple operations per second, and about $64$ times more if they are word operations.

\section*{Step 1 --- What the limits allow}

\begingroup\small\renewcommand{\arraystretch}{1.3}
\begin{tabularx}{\linewidth}{@{}l l X@{}}
\toprule
\textbf{Limit} & \textbf{Budget} & \textbf{Typical bit techniques} \\
\midrule
$n \le 20$ & $\bigO{2^n n}$ & enumerate all subsets as masks (Ch.~\ref{ch:subsets}); bitmask DP is the DP book \\
$n \le 15$--$16$ & $\bigO{3^n}$ & all submasks of all masks (Ch.~\ref{ch:mask_enum}) \\
$n \le 40$ & $\bigO{2^{n/2}}$ & meet in the middle (the \textit{Recursion} book) \\
$V \le 2^{20}$ & $\bigO{V \log V}$ & Walsh--Hadamard and AND/OR transforms over the value domain (Ch.~\ref{ch:fwht}) \\
$n \le 2000$--$5000$ & $\bigO{n^3/64}$, $\bigO{n^2}$ & bitsets: Gaussian elimination over GF(2), transitive closure (Ch.~\ref{ch:gf2}, \ref{ch:bitsets}) \\
$n \le 2\cdot10^5$, $V \le 2^{30}$ & $\bigO{n \log V}$ & per-bit counting, greedy over bits, linear basis insertions (Ch.~\ref{ch:per_bit}, \ref{ch:greedy_bits}, \ref{ch:linear_basis}) \\
$n \le 10^6$ & $\bigO{n}$ & prefix XOR, popcount per element, XOR hashing (Ch.~\ref{ch:popcount}, \ref{ch:prefix_xor}) \\
$V$ up to $10^{18}$, one number & $\bigO{\log V}$ & walk its bits: closed forms, counting over binary prefixes (Ch.~\ref{ch:xor}, \ref{ch:binary_prefix}) \\
\bottomrule
\end{tabularx}
\endgroup

\begin{insightbox}[Read the width, not only the value]
$V \le 10^9$ fits in \texttt{int}, but $1 \ll 31$ does not; $V \le 10^{18}$ needs \texttt{long long} and every shift must be \texttt{1LL << k}. A limit of $2^{60}$ means a linear basis has at most $60$ vectors and a greedy over bits runs $60$ rounds.
\end{insightbox}

\clearpage
\section*{Step 2 --- What the question asks}

\begingroup\small\renewcommand{\arraystretch}{1.3}
\begin{tabularx}{\linewidth}{@{}>{\raggedright\arraybackslash}p{6.2cm} X@{}}
\toprule
\textbf{The problem says \ldots} & \textbf{Try first} \\
\midrule
``binary representation'', ``number of 1 bits'', ``complement'' & read the bits (Ch.~\ref{ch:introduction}, \ref{ch:popcount}) \\
``power of two'', ``lowest set bit'', ``trailing zeros'' & classic tricks (Ch.~\ref{ch:tricks}) \\
XOR of a range of integers, decoding XORed arrays & XOR closed forms and cancellation (Ch.~\ref{ch:xor}) \\
``all subsets'', $n \le 20$ & masks as subsets (Ch.~\ref{ch:subsets}) \\
``maximum XOR / AND of a pair'' & greedy from the highest bit (Ch.~\ref{ch:greedy_bits}) \\
``sum of XOR / AND / OR over all pairs'' & per-bit contribution (Ch.~\ref{ch:per_bit}) \\
given $a + b$ and $a \oplus b$ (or OR and AND) & bitwise equations (Ch.~\ref{ch:equations}) \\
``flip one bit at a time'', ``adjacent differ in one bit'' & Gray codes and hypercubes (Ch.~\ref{ch:gray}) \\
``maximum XOR of a subset'', ``how many subsets have XOR $x$'' & linear basis (Ch.~\ref{ch:linear_basis}) \\
``toggle'' puzzles, systems of XOR equations & Gaussian elimination over GF(2) (Ch.~\ref{ch:gf2}) \\
``count pairs with XOR $= k$'' over a value domain & Walsh--Hadamard transform (Ch.~\ref{ch:fwht}) \\
two players, piles, take any amount & Nim and XOR (Ch.~\ref{ch:nim}) \\
``count $x \le n$ such that \ldots'' with a bitwise condition & counting over binary prefixes (Ch.~\ref{ch:binary_prefix}) \\
hidden numbers, ``ask AND / OR / XOR queries'' & interactive bit problems (Ch.~\ref{ch:interactive}) \\
``construct an array with XOR \ldots'' & constructive bit problems (Ch.~\ref{ch:constructive}) \\
XOR / AND / OR over all subarrays & the \textit{Arrays} book \\
maximum XOR with a set under insertions and deletions & binary tries (the \textit{Strings} and \textit{Trees} books) \\
DP over subsets, SOS, digit DP & the \textit{Dynamic Programming} book \\
\bottomrule
\end{tabularx}
\endgroup

\begin{warningbox}[Limits are hints, not proofs]
The table gives the \emph{intended} complexity in most problems, not a guarantee. A small $n$ with $V \le 2^{60}$ may hide a linear basis instead of subset enumeration; a ``sum over all pairs'' with $n \le 2\cdot10^5$ is always per-bit, never pairwise.
\end{warningbox}
